% !TeX root = ../../../geos_mpm_manual.tex
\section{Boundary conditions, prescribed deformation, and stress control}
\label{sec:bc-internals}
\index{boundary conditions}
\index{Ftable}
\index{stress control}
\index{reactionHistory}

This section describes what the MPM solver does internally when boundary and domain-control options are active.  It deliberately separates the numerical implementation from the user-facing input syntax.  The user-facing \texttt{pfw\_input} controls, including examples of \texttt{boundaryConditionTypes}, \texttt{fTable}, \texttt{bcTable}, transverse boundary velocities, stress control, and periodic flags, are summarized in Section~\ref{sec:pfw-boundary-controls}.  The code path described here is the one reached during Step~14 of the explicit solver update, after particle-to-grid projection, MPI synchronization, grid trial dynamics, and optional material-field contact; see Section~\ref{subsec:solver-step-14}.

\subsection{Face convention and nodal data used by the boundary update}
\label{subsec:bc-face-convention}
\index{boundary conditions!face order}

The solver stores one boundary type for each axis-aligned face of the background grid.  The face order is
\begin{equation}
  f=0,1,2,3,4,5
  \quad\leftrightarrow\quad
  x^-,x^+,y^-,y^+,z^-,z^+.
\end{equation}
For face $f$, the normal coordinate direction is
\begin{equation}
  d(f)=\left\lfloor \frac{f}{2}\right\rfloor,
  \qquad
  s(f)=f\bmod 2,
  \qquad
  \alpha_f=-1+2s(f),
  \qquad
  \mathbf{n}_f=\alpha_f\mathbf{e}_{d(f)}.
  \label{eq:bc-face-normal}
\end{equation}
The two remaining coordinate directions are used as tangential directions.  In the code these are computed as
\begin{equation}
  d_1=(d+1)\bmod 3,
  \qquad
  d_2=(d+2)\bmod 3.
\end{equation}

Boundary enforcement operates on grid multi-fields.  The field index distinguishes material/contact velocity fields, while the node index distinguishes ordinary grid nodes, boundary nodes, and buffer/ghost-like nodes used to support the interpolation stencil near a face.  The principal fields modified by the boundary code are the grid velocity $\mathbf{v}_{I\alpha}$, grid acceleration $\mathbf{a}_{I\alpha}$, and grid velocity increment $\Delta\mathbf{v}_{I\alpha}$ for node $I$ and velocity field $\alpha$.  The boundary update also reads nodal mass, internal force, center-of-volume, surface-position, surface-field-mass, and ghost-rank data when contact or reaction histories are active.

The boundary-node and buffer-node sets are stored in the node-set repository as \texttt{boundaryNodes} and \texttt{bufferNodes}.  For a face $f$, let $\mathcal{B}_f$ be its boundary nodes and $\mathcal{G}_f$ its buffer nodes.  Each buffer node $G\in\mathcal{G}_f$ is paired with an interior node $J(G)$ and, for moving faces, a boundary node $B(G)$ at the same tangential grid coordinates.  These pairings are reconstructed from the structured-grid \texttt{ijkMap} and the current grid spacing.  They provide the even/odd extensions needed by B-spline and other wide-stencil particle-grid transfers near a boundary.

\subsection{Boundary-type state and time-dependent boundary-type switching}
\label{subsec:bc-table-internals}
\index{boundary conditions!bcTable}
\index{bcTable}

The runtime boundary-type array is \texttt{m\_boundaryConditionTypes}.  If the input does not provide this array, it is initialized to six \texttt{Outflow} entries.  If it is provided, the solver requires exactly six integer entries and checks that each entry belongs to the boundary-condition enumeration in Table~\ref{tab:bc-types}.

\begin{table}[htbp]
\centering
\caption{MPM boundary-condition type enumeration used by \texttt{SolidMechanicsMPM}.}
\label{tab:bc-types}
\begin{tabularx}{\linewidth}{@{}c l X@{}}
\toprule
Integer & Name & Internal interpretation \\
\midrule
0 & \texttt{Outflow} & Leave grid velocity and acceleration unconstrained on that face. \\
1 & \texttt{Symmetry} & Enforce zero normal velocity/acceleration and reflect buffer fields with an odd normal component and even tangential components. \\
2 & \texttt{Moving} & Enforce a prescribed affine normal velocity when prescribed boundary deformation or stress control is active; optionally enforce tangential velocities. \\
3 & \texttt{Contact} & Treat the face as a wall/contact plane.  Incoming nodes are corrected in the face-normal direction and optional boundary friction is applied in the tangential plane. \\
\bottomrule
\end{tabularx}
\end{table}

If \texttt{prescribedBcTable} is enabled, Step~4 of the explicit update calls \texttt{boundaryConditionUpdate} before particle-grid maps are assembled.  The table is a zero-order-hold schedule rather than an interpolated quantity.  For a time step beginning at $t^n$ with size $\Delta t$, the selected row is the last row whose time satisfies
\begin{equation}
  t_{\mathrm{row}} < t^n+\frac{\Delta t}{2}.
  \label{eq:bc-table-selection}
\end{equation}
The six boundary-type entries in that row then replace \texttt{m\_boundaryConditionTypes}.  After this update, \texttt{updateContactFlagsFromBoundaryConditions} sets \texttt{m\_hasContact} if either multiple velocity fields are present or at least one face is marked \texttt{Contact}.  This flag controls whether contact/surface fields are synchronized and whether contact-specific work is performed later in the step.

\begin{lstlisting}[caption={Boundary-type update and contact flag refresh.}]
if prescribedBcTable == 1:
  interval = last row i such that time_n + 0.5*dt > bcTable[i,0]
  for face in 0..5:
    boundaryConditionTypes[face] = bcTable[interval, face+1]

hasContact = (numVelocityFields > 1) or any(boundaryConditionTypes == Contact)
\end{lstlisting}

\subsection{Outflow faces}
\label{subsec:bc-outflow}
\index{boundary conditions!Outflow}

An \texttt{Outflow} face is the least intrusive boundary option.  In \texttt{applyEssentialBCs}, the face is skipped:
\begin{equation}
  \mathbf{v}_{I\alpha}^{\mathrm{after}}=\mathbf{v}_{I\alpha}^{\mathrm{before}},
  \qquad
  \mathbf{a}_{I\alpha}^{\mathrm{after}}=\mathbf{a}_{I\alpha}^{\mathrm{before}},
  \qquad I\in\mathcal{B}_f.
\end{equation}
No reflective buffer update is applied by the essential-boundary routine for this face.  This should not be interpreted as a finite-element natural traction boundary.  It means that the explicit MPM grid solve leaves the nodal values produced by particle-to-grid mapping, force assembly, contact, and MPI synchronization unchanged.  Particles that advect beyond a nonperiodic global domain are handled later by the out-of-range particle flagging and cleanup path.  Consequently, \texttt{Outflow} is useful for open-domain style calculations, impact problems where material may leave the computational box, or inactive directions in which the user does not want a symmetry or prescribed-motion constraint.

\subsection{Symmetry and free-slip reflection}
\label{subsec:bc-symmetry}
\index{boundary conditions!Symmetry}
\index{symmetry boundary conditions}

A \texttt{Symmetry} face enforces a kinematic free-slip plane.  On the boundary nodes, the normal component of the vector field is removed.  For velocity and acceleration this gives
\begin{equation}
  v_{I\alpha,d}=0,
  \qquad
  a_{I\alpha,d}=0,
  \qquad I\in\mathcal{B}_f.
  \label{eq:bc-symmetry-boundary-node}
\end{equation}
The corresponding velocity increment is updated so that XPIC/FMPM-style transfer formulas see the boundary-induced correction,
\begin{equation}
  \Delta v_{I\alpha,d} \leftarrow \Delta v_{I\alpha,d} - v_{I\alpha,d}^{\mathrm{old}}.
\end{equation}

For buffer nodes, the solver uses an odd extension of the normal component and an even extension of the tangential components from the interior paired node $J(G)$:
\begin{align}
  v_{G\alpha,d}   &= -v_{J(G)\alpha,d}, &
  v_{G\alpha,d_1} &=  v_{J(G)\alpha,d_1}, &
  v_{G\alpha,d_2} &=  v_{J(G)\alpha,d_2},
  \label{eq:bc-symmetry-buffer-velocity} \\
  a_{G\alpha,d}   &= -a_{J(G)\alpha,d}, &
  a_{G\alpha,d_1} &=  a_{J(G)\alpha,d_1}, &
  a_{G\alpha,d_2} &=  a_{J(G)\alpha,d_2}.
\end{align}
This mirrored extension is also applied earlier to mapped surface normals, center-of-mass vectors, and center-of-volume vectors on faces marked \texttt{Symmetry} or \texttt{Moving}.  That earlier pass occurs before the grid dynamics update and ensures that contact and surface quantities remain compatible with the same geometric reflection used for the final velocity and acceleration fields.

The construction is the MPM analogue of imposing an essential normal-velocity condition on the background grid while allowing tangential slip.  It is consistent with the grid-based enforcement commonly used in explicit MPM and GIMP calculations, where the particle state is advanced using constrained nodal accelerations and velocities rather than by imposing constraints directly on the particles \cite{sulsky1994history,bardenhagen2004gimp,wallstedt2008explicit}.

\begin{lstlisting}[caption={Symmetry-face enforcement.}]
for field alpha:
  for I in boundaryNodes[face]:
    dVelocity[I,alpha,d] = -velocity[I,alpha,d]
    velocity[I,alpha,d] = 0
    acceleration[I,alpha,d] = 0

  for G in bufferNodes[face]:
    J = reflected interior node for G
    velocity[G,alpha,d]  = -velocity[J,alpha,d]
    velocity[G,alpha,t1] =  velocity[J,alpha,t1]
    velocity[G,alpha,t2] =  velocity[J,alpha,t2]
    acceleration[G,alpha,d]  = -acceleration[J,alpha,d]
    acceleration[G,alpha,t1] =  acceleration[J,alpha,t1]
    acceleration[G,alpha,t2] =  acceleration[J,alpha,t2]
\end{lstlisting}

\subsection{Prescribed domain deformation and F-table interpolation}
\label{subsec:bc-ftable-internals}
\index{Ftable!interpolation}
\index{prescribedBoundaryFTable}
\index{prescribedFTable}

The F-table machinery produces a diagonal macroscopic deformation history and its associated diagonal velocity-gradient components.  The table has one time column and three deformation columns.  At the end of the step, $t^{n+1}=t^n+\Delta t$, \texttt{interpolateTable} selects the interval containing the actual evaluation time $t^{n+1}$ and returns $F_i(t^{n+1})$ and $\dot F_i(t^{n+1})$.  This endpoint selection avoids extrapolating from the interval containing the time-step midpoint when a step crosses a table knot.  At an internal knot the derivative is taken from the interval on the left; before the first time and after the final time the value is clamped and the rate is zero.  For directions not completely overridden by stress control, the code sets
\begin{equation}
  \bar L_i(t^{n+1}) = \frac{\dot F_i(t^{n+1})}{F_i(t^{n+1})},
  \qquad
  \bar F_i(t^{n+1}) = F_i(t^{n+1}).
  \label{eq:bc-ftable-l}
\end{equation}
For \texttt{prescribedBoundaryFTable}, the solver separately forms the finite grid map
\begin{equation}
  \lambda_i^{n\rightarrow n+1}
  = \frac{\bar F_i(t^{n+1})}{\bar F_i(t^n)}.
  \label{eq:bc-ftable-incremental-stretch}
\end{equation}
The distinction is intentional: $\bar L_i(t^{n+1})$ is the instantaneous endpoint rate used for wall velocity, whereas $\lambda_i^{n\rightarrow n+1}$ is the exact finite stretch used for grid coordinates and domain geometry.
The current implementation supports three scalar interpolation kernels.  Let $y_0,y_1$ be consecutive table values, $\tau=(t^{n+1}-t_0)/(t_1-t_0)$, and $\Delta T=t_1-t_0$.  Linear interpolation uses
\begin{equation}
  y(\tau)=(1-\tau)y_0+\tau y_1,
  \qquad
  \dot y=\frac{y_1-y_0}{\Delta T}.
\end{equation}
Cosine interpolation uses
\begin{equation}
  y(\tau)=y_0-\frac{1}{2}(y_1-y_0)\left[\cos(\pi\tau)-1\right],
  \qquad
  \dot y=\frac{\pi}{2\Delta T}(y_1-y_0)\sin(\pi\tau),
\end{equation}
and \texttt{Smoothstep} uses the quintic polynomial
\begin{equation}
  y(\tau)=y_0+(y_1-y_0)(10\tau^3-15\tau^4+6\tau^5),
  \qquad
  \dot y=\frac{y_1-y_0}{\Delta T}(30\tau^2-60\tau^3+30\tau^4).
\end{equation}
Values before the first table time are clamped to the first row, and values after the final table time are clamped to the last row.

The two F-table flags route this macroscopic kinematics into different parts of the MPM update.  With \texttt{prescribedBoundaryFTable}, the diagonal endpoint velocity-gradient components become boundary velocities on \texttt{Moving} faces during \texttt{applyEssentialBCs}, while Eq.~\eqref{eq:bc-ftable-incremental-stretch} updates wall positions, grid nodes, partition/domain extents, reaction-history dimensions, and rigid-body jamming areas.  With \texttt{prescribedFTable}, the rate components are superposed on particle velocity gradients and positions only in directions marked periodic by the spatial partition.  The latter path is intended for triply periodic or partially periodic homogeneous-deformation calculations and retains its existing rate-based update.

\subsection{Moving faces}
\label{subsec:bc-moving}
\index{boundary conditions!Moving}
\index{prescribedBoundaryTransverseVelocities}

A \texttt{Moving} face is an essential velocity boundary driven by the current domain velocity-gradient state.  In the reviewed implementation, the velocity constraint is applied only if either \texttt{prescribedBoundaryFTable} is active or at least one \texttt{stressControl} entry equals one.  This guard is important: marking a face as \texttt{Moving} without an active prescribed boundary deformation, or with only the limiting \texttt{stressControl} mode two active, does not by itself prescribe a nonzero velocity in \texttt{applyEssentialBCs}.

For a moving face $f$ normal to direction $d$, the code forms an end-of-step grid coordinate
\begin{equation}
  X_{I,d}^{n+1,*}=\lambda_d^{n\rightarrow n+1}X_{I,d}^{n}.
  \label{eq:bc-moving-exact-position}
\end{equation}
For \texttt{prescribedBoundaryFTable}, this is the exact tabulated position because $\lambda_d^{n\rightarrow n+1}=\bar F_d^{n+1}/\bar F_d^n$.  For stress-control-only and periodic-deformation paths, the stored factor retains the existing $1+\bar L_d\Delta t$ form.  The normal velocity is
\begin{equation}
  v_{I\alpha,d}^{\mathrm{bc}} = \bar L_d(t^{n+1}) X_{I,d}^{n+1,*},
  \qquad I\in\mathcal{B}_f.
  \label{eq:bc-moving-normal-velocity}
\end{equation}
The solver then computes the increment and acceleration correction as
\begin{equation}
  \Delta v_{I\alpha,d} \leftarrow v_{I\alpha,d}^{\mathrm{bc}}-v_{I\alpha,d}^{\mathrm{old}},
  \qquad
  a_{I\alpha,d} \leftarrow a_{I\alpha,d}+\frac{v_{I\alpha,d}^{\mathrm{bc}}-v_{I\alpha,d}^{\mathrm{old}}}{\Delta t},
  \qquad
  v_{I\alpha,d}\leftarrow v_{I\alpha,d}^{\mathrm{bc}}.
  \label{eq:bc-moving-normal-correction}
\end{equation}

Tangential motion is optional.  If \texttt{enablePrescribedBoundaryTransverseVelocities[f]} is one, the two stored tangential values are enforced on $d_1$ and $d_2$ exactly as essential velocity components:
\begin{equation}
  v_{I\alpha,d_1}\leftarrow v_{f,d_1}^{\mathrm{presc}},
  \qquad
  v_{I\alpha,d_2}\leftarrow v_{f,d_2}^{\mathrm{presc}},
  \label{eq:bc-moving-transverse}
\end{equation}
with matching updates to \texttt{gridDVelocity} and \texttt{gridAcceleration}.  If transverse velocities are not enabled, the tangential components are left free on the boundary nodes.

The buffer-node extension is centered on the prescribed boundary value.  For the constrained normal component,
\begin{equation}
  v_{G\alpha,d}=-v_{J(G)\alpha,d}+2v_{B(G)\alpha,d}.
  \label{eq:bc-moving-buffer-normal}
\end{equation}
For tangential components, the code applies the same moving odd extension only when transverse velocities are explicitly prescribed.  When transverse motion is not prescribed, the buffer-node tangential velocity, increment, and acceleration are left unchanged:
\begin{equation}
  v_{G\alpha,d_k}=\begin{cases}
    -v_{J(G)\alpha,d_k}+2v_{B(G)\alpha,d_k}, & \text{prescribed transverse component active},\\
    v_{G\alpha,d_k}^{\mathrm{before}}, & \text{transverse component free},
  \end{cases}
  \qquad k=1,2.
  \label{eq:bc-moving-buffer-tangent}
\end{equation}
This construction constrains only the physical Dirichlet components.  Free tangential components are natural/slip degrees of freedom; overwriting them on mass-bearing buffer/control nodes can inject a real tangential impulse.

\begin{lstlisting}[caption={Moving-face enforcement.}]
for field alpha and boundary node I:
  x_end = domainIncrementalStretch[d] * gridPosition[I,d]
  v_bc = domainL[d] * x_end

  if transverse velocities enabled on this face:
    set velocity[I,alpha,t1] and velocity[I,alpha,t2]
    update dVelocity and acceleration corrections in t1,t2

  dVelocity[I,alpha,d] = v_bc - velocity[I,alpha,d]
  acceleration[I,alpha,d] += (v_bc - velocity[I,alpha,d]) / dt
  velocity[I,alpha,d] = v_bc

for buffer node G:
  J = reflected interior node
  B = boundary node at same tangential indices
  velocity[G,d] = -velocity[J,d] + 2*velocity[B,d]
  if tangential constrained:
    velocity[G,t] = -velocity[J,t] + 2*velocity[B,t]
  else:
    leave velocity[G,t] unchanged
\end{lstlisting}

\subsection{Boundary contact faces}
\label{subsec:bc-contact-wall}
\index{boundary conditions!Contact}
\index{boundary friction}

A face marked \texttt{Contact} is handled in the same \texttt{Moving}/\texttt{Contact} branch, but with wall-contact detection replacing unconditional prescribed motion.  The branch is always entered for a contact face, even if no F-table or stress-control flag is active.  For each boundary node and velocity field, the code first estimates a scalar surface position in the face-normal direction.  If a particle-mapped surface position is available at the node, it uses \texttt{gridSurfacePosition}; otherwise it falls back to \texttt{gridCenterOfVolume}.  The predicted wall velocity is
\begin{equation}
  v_{I,d}^{\mathrm{wall}} = \bar L_d(t^{n+1})\lambda_d^{n\rightarrow n+1}X_{I,d}^{n},
  \label{eq:bc-wall-normal-target}
\end{equation}
and the relative normal velocity is $v_{I\alpha,d}^{\mathrm{rel}}=v_{I\alpha,d}-v_{I,d}^{\mathrm{wall}}$.  With the sign $\alpha_f$ from Eq.~\eqref{eq:bc-face-normal}, the current and predicted outward surface coordinates are
\begin{equation}
  s^n=\alpha_f x^{\mathrm{surf}}_{I\alpha,d},
  \qquad
  s^{n+1,*}=\alpha_f\left(x^{\mathrm{surf}}_{I\alpha,d}+\Delta t\,v_{I\alpha,d}^{\mathrm{rel}}\right).
\end{equation}
The normal constraint activates when the surface intersects the wall, $\max(s^n,s^{n+1,*})>0$, and the material moves into it, $\alpha_f v_{I\alpha,d}^{\mathrm{rel}}>0$; the rigid-body penetration penalty can also activate an intersecting wall.  When contact is inactive, the face does not overwrite the normal velocity.  When it is active, the normal boundary target is $v_{I\alpha,d}^{\mathrm{bc}}=v_{I,d}^{\mathrm{wall}}$.
followed by the same normal velocity-increment, acceleration-correction, reaction-accumulation, and buffer-reflection operations used by a moving face.  The registered \texttt{boundaryFaceCoefficientsOfRestitution} array is initialized and checked, but the active code path in the reviewed source does not use it in the normal target; a restitution-based line is present only as commented legacy code.  Thus boundary contact should currently be interpreted as a sticking/non-penetration style normal constraint with optional tangential friction, not as a coefficient-of-restitution bounce law.

Tangential friction is controlled by \texttt{boundaryFaceFrictionCoefficients}.  Let
\begin{equation}
  \mathbf{v}_t = v_{I\alpha,d_1}\mathbf{e}_{d_1}+v_{I\alpha,d_2}\mathbf{e}_{d_2},
  \qquad
  \mathbf{r}_t = \begin{cases}
    \mathbf{v}_t/\lVert\mathbf{v}_t\rVert, & \lVert\mathbf{v}_t\rVert>10^{-16},\\
    \mathbf{0}, & \text{otherwise},
  \end{cases}
\end{equation}
and define the code's scalar friction measure by
\begin{equation}
  f_\mu = \mu_f\max\left(-\alpha_f a_{I\alpha,d},0\right).
  \label{eq:bc-wall-friction-scalar}
\end{equation}
The code then updates the tangential velocity increment and acceleration fields component-wise as
\begin{equation}
  \Delta\mathbf{v}_t \leftarrow f_\mu\Delta t\,\mathbf{r}_t,
  \qquad
  \mathbf{v}_t \leftarrow \mathbf{v}_t+\Delta\mathbf{v}_t,
  \qquad
  \mathbf{a}_t \leftarrow \mathbf{a}_t-f_\mu\mathbf{r}_t.
  \label{eq:bc-wall-friction-update}
\end{equation}
Equation~\eqref{eq:bc-wall-friction-update} is intentionally written to match the reviewed implementation.  Because the sign convention combines face orientation, acceleration direction, and the subsequent velocity/acceleration updates, wall-friction verification cases should be used when relying on nonzero boundary friction.

\begin{lstlisting}[caption={Contact-face wall update.}]
for field alpha and boundary node I:
  surfacePosition = gridSurfacePosition if available else gridCenterOfVolume
  incoming = alpha(face)*velocity[I,alpha,d] > 0
  outside  = alpha(face)*surfacePosition[d] > 0

  if incoming and outside:
    v_bc = domainL[d] * gridPosition[I,d]
    apply normal velocity correction as for Moving

  mu = boundaryFaceFrictionCoefficients[face]
  frictionMeasure = mu * max(-alpha(face)*acceleration[I,alpha,d], 0)
  apply implemented tangential velocity/acceleration update
\end{lstlisting}

\subsection{Stress-control generated boundary motion}
\label{subsec:bc-stress-control-internals}
\index{stress control!internal algorithm}
\index{stressTable}

Stress control does not impose tractions directly on grid faces.  Instead, it computes a diagonal domain velocity gradient, stores it in \texttt{domainL}, and then lets the moving-boundary machinery impose the associated affine boundary velocities.  The target stress is obtained from \texttt{stressTable} using the same interpolation framework as the F-table.  The current stress is computed by \texttt{computeBoxMetrics} from particle stresses and material volume in the solver's averaging box.  In compact notation,
\begin{equation}
  e_i = \sigma_i^{\mathrm{target}}-\sigma_i^{\mathrm{current}}.
  \label{eq:bc-stress-control-error}
\end{equation}

The controller scales the input gains by an estimate of relative density and the maximum bulk or effective bulk modulus over active particles.  If $K_{\max}$ is the global maximum bulk modulus and
\begin{equation}
  \rho_{\mathrm{rel}}=\min\left(1,\frac{V_{\mathrm{mat}}}{V_{\mathrm{domain}}}\right),
\end{equation}
then the effective gains used by the code are
\begin{equation}
  k_p=\frac{\rho_{\mathrm{rel}}\,\texttt{stressControlKp}}{K_{\max}\Delta t},
  \qquad
  k_d=\frac{\rho_{\mathrm{rel}}\,\texttt{stressControlKd}}{K_{\max}},
  \qquad
  k_i=\frac{\rho_{\mathrm{rel}}\,\texttt{stressControlKi}}{K_{\max}\Delta t^2}.
  \label{eq:bc-stress-control-scaled-gains}
\end{equation}
The integral state and derivative estimate are updated as
\begin{align}
  \mathbf{I}^{n+1} &= \mathbf{I}^{n}+k_i\Delta t\,\mathbf{e}^{n+1},\\
  \dot{\mathbf{e}}^{n+1} &= \frac{\mathbf{e}^{n+1}-\mathbf{e}^{n}}{\Delta t},
\end{align}
and the unconstrained controller output is
\begin{equation}
  \bar{\mathbf{L}}_{\mathrm{PID}} = k_p\mathbf{e}^{n+1}+\mathbf{I}^{n+1}+k_d\dot{\mathbf{e}}^{n+1}.
  \label{eq:bc-stress-control-pid}
\end{equation}
Each component is clipped so that the implied boundary speed is bounded by the CFL length scale,
\begin{equation}
  |\bar L_i| \le \frac{\texttt{cflFactor}\,h_i}{\Delta t\,(X_i^{\max}-X_i^{\min})}.
  \label{eq:bc-stress-control-cfl-limit}
\end{equation}
For \texttt{stressControl[i] == 1}, the clipped value replaces \texttt{domainL[i]}.  For \texttt{stressControl[i] == 2}, the direction is treated as a limiting mode that prevents further deformation in the forbidden direction once the stress error changes sign.  The updated \texttt{domainL} is then converted to \texttt{domainF} by the explicit update
\begin{equation}
  \bar F_i^{n+1}=\bar F_i^n+\bar L_i\bar F_i^n\Delta t.
\end{equation}
Moving faces normal to controlled directions then use Eq.~\eqref{eq:bc-moving-normal-velocity} when the moving-face branch is active; in practice, stress-controlled boundary-motion cases should use \texttt{stressControl[i] == 1} for controlled directions and commonly retain \texttt{prescribedBoundaryFTable} with a neutral F-table for robust activation.  In a stress-controlled direction, the explicit update above gives $\bar F_i^{n+1}/\bar F_i^n=1+\bar L_i\Delta t$, so the finite-ratio geometry path reproduces the previous controller motion while table-controlled directions obtain the exact prescribed ratio.  In this sense, stress control is internally a feedback generator for prescribed kinematics rather than a separate traction-boundary algorithm.

\begin{lstlisting}[caption={Stress-control path to boundary motion.}]
interpolateStressTable(time_n, dt) -> domainStress
computeBoxMetrics() -> currentStress, materialVolume
Kmax = global maximum bulk/effective bulk modulus over active particles
relativeDensity = min(1, materialVolume / domainVolume)
error = targetStress - currentStress
L_pid = scaled_Kp*error + integral + scaled_Kd*(error-lastError)/dt
clip L_pid by CFL boundary speed
for each controlled direction:
  domainL[i] = L_pid[i] or limiter-modified value
  domainF[i] = oldDomainF[i] + domainL[i]*oldDomainF[i]*dt
apply moving-face constraints using domainL
\end{lstlisting}

\subsection{Adaptive-seek normal stress control}
\label{subsec:bc-adaptive-seek-stress-control}
\index{stress control!adaptive-seek}
\index{stressControlAdaptiveSeek}

The optional adaptive-seek controller is enabled by setting
\texttt{stressControlAdaptiveSeek = 1}.  It replaces the legacy PID update in
Subsection~\ref{subsec:bc-stress-control-internals} while preserving the same
\texttt{stressControl} and \texttt{stressTable} interface.  The current
implementation controls only the diagonal normal components represented by
\texttt{domainL} and \texttt{domainF}; shear stress control and non-orthogonal
periodic-cell control are outside its current scope.  In adaptive-seek mode,
\texttt{stressControl[i]} entries must be either zero or one.  The legacy
one-sided limiter value two is not supported by this controller.

The adaptive controller is intended for mixed prescribed-F and stress-controlled
RVE calculations where a scalar bulk-modulus normalization is a poor
approximation.  It learns a small normal tangent,
\begin{equation}
  \Delta\sigma_i \approx \sum_{j=0}^{2} C_{ij}\Delta\epsilon_j,
  \qquad i,j\in\{xx,yy,zz\},
  \label{eq:bc-adaptive-seek-tangent}
\end{equation}
from filtered box-average stresses and the applied normal velocity-gradient
history.  The diagonal tangent is initialized from the same maximum bulk or
effective bulk modulus used by the legacy controller, multiplied by the relative
material density; this is only an initial scale and is updated during loading.

Let $A$ be the set of stress-controlled normal components and $P$ the remaining
normal components.  The active controller computes a damped least-squares solve
of the form
\begin{equation}
  \mathbf{L}_A = C_{AA}^{+}\left(
    -C_{AP}\mathbf{L}_P
    + \frac{\boldsymbol{\sigma}^{\mathrm{target}}_A-
             \boldsymbol{\sigma}^{\mathrm{filtered}}_A}
           {\epsilon_{\mathrm{response}}}
      L_{\mathrm{ref}}
  \right),
  \label{eq:bc-adaptive-seek-active-law}
\end{equation}
where $C_{AA}^{+}$ denotes the damped small-system solve, not a raw inverse.
The first term provides feed-forward coupling from prescribed F-table
directions, and the second term removes filtered stress error over the strain
interval \texttt{stressControlResponseStrain}.  The reference rate
$L_{\mathrm{ref}}$ is the maximum of the prescribed/inactive normal rates, the
previous controlled rates, and \texttt{stressControlMinStrainRate}.  Thus the
controller bandwidth follows the macroscopic deformation path rather than the
explicit timestep or an assumed wave speed.

The raw stress is low-pass filtered over a strain interval,
\begin{equation}
  \alpha = \min\left(1,
    \frac{L_{\mathrm{ref}}\Delta t}{\texttt{stressControlFilterStrain}}
  \right),
  \qquad
  \sigma_i^{\mathrm{filtered}} \leftarrow
  \sigma_i^{\mathrm{filtered}} +
  \alpha(\sigma_i^{\mathrm{raw}}-\sigma_i^{\mathrm{filtered}}).
  \label{eq:bc-adaptive-seek-filter}
\end{equation}
The largest normalized raw-minus-filtered stress difference over active
components defines a wave or inertia indicator.  If it exceeds
\texttt{stressControlWaveWarnRatio}, the controller prints a warning, slows the
feedback response, and freezes tangent adaptation while the condition persists.
This diagnostic is designed to detect cases where prescribed F-table loading is
generating stress waves that dominate the measured box stress.

If a compressive target has no current load path, a zero tangent should not be
inverted.  The adaptive-seek controller therefore switches to bounded seek mode
when the target pressure $p_i^{\mathrm{target}}=-\sigma_i^{\mathrm{target}}$ is
compressive and the measured filtered pressure is too small.  In this state,
\begin{equation}
  L_i^{\mathrm{seek}} =
  -L_{\mathrm{seek}}
  \operatorname{clip}\left(
    \frac{p_i^{\mathrm{target}}-p_i^{\mathrm{filtered}}}
         {\max(|p_i^{\mathrm{target}}|,p_{\mathrm{floor}})},0,1
  \right),
  \label{eq:bc-adaptive-seek-seek-law}
\end{equation}
where $L_{\mathrm{seek}}$ is limited by \texttt{stressControlSeekRateRatio} and
\texttt{stressControlMaxSeekRateRatio}.  Seek mode re-engages active control
after a pressure fraction \texttt{stressControlJammingPressureRatio} appears and
the wave indicator is below \texttt{stressControlWaveCutbackRatio}.  The return
to active control is blended over \texttt{stressControlReengageRampStrain}.  If
\texttt{stressControlMaxSeekStrain} is exceeded in a direction, that direction is
marked unreachable and its controlled rate is held at zero.

\begin{table}[h]
\centering
\caption{Main adaptive-seek stress-control parameters.}
\label{tab:bc-adaptive-seek-parameters}
\begin{tabular}{lll}
\toprule
Parameter & Default & Purpose \\
\midrule
\texttt{stressControlAdaptiveSeek} & 0 & Enable adaptive-seek control. \\
\texttt{stressControlResponseStrain} & $10^{-2}$ & Strain interval for stress-error correction. \\
\texttt{stressControlFilterStrain} & $10^{-3}$ & Strain interval for stress filtering. \\
\texttt{stressControlAdaptStrainWindow} & $5\times10^{-4}$ & Strain window for tangent updates. \\
\texttt{stressControlAdaptGain} & $5\times10^{-2}$ & Rank-one tangent update gain. \\
\texttt{stressControlMaxRateRatio} & 5 & Total active rate cap relative to reference rate. \\
\texttt{stressControlMaxFeedbackRateRatio} & 2 & Feedback-only rate cap. \\
\texttt{stressControlSeekRateRatio} & 0.5 & Bounded seek rate ratio. \\
\texttt{stressControlMaxSeekRateRatio} & 2 & Maximum seek rate ratio. \\
\texttt{stressControlMinStrainRate} & $10^{-30}$ & Minimum reference rate. \\
\texttt{stressControlMaxSeekStrain} & $5\times10^{-2}$ & Seek strain before declaring unreachable. \\
\texttt{stressControlMaxSeekStrainIncrement} & $10^{-4}$ & Per-step strain increment cap. \\
\texttt{stressControlJammingPressureRatio} & $2\times10^{-2}$ & Pressure fraction for load-path appearance. \\
\texttt{stressControlReengageRampStrain} & $5\times10^{-3}$ & Seek-to-active blending strain. \\
\texttt{stressControlWaveWarnRatio} & $5\times10^{-2}$ & Raw-filtered stress warning ratio. \\
\texttt{stressControlWaveCutbackRatio} & $10^{-1}$ & Re-engagement/adaptation cutback ratio. \\
\texttt{stressControlStressFloor} & 0 & Explicit stress floor; target-based if zero. \\
\texttt{stressControlTangentFloorRatio} & $10^{-3}$ & Lower tangent bound ratio. \\
\texttt{stressControlTangentCeilingRatio} & $10^{2}$ & Upper tangent bound ratio. \\
\texttt{stressControlMaxCouplingRatio} & 10 & Off-diagonal tangent cap. \\
\texttt{stressControlSolverDampingRatio} & $5\times10^{-2}$ & Damped-solve regularization. \\
\bottomrule
\end{tabular}
\end{table}

For prescribed-F quasi-static tests, the default strain scales are suitable
starting values.  For multiple controlled directions, use a smaller tangent
adaptation gain, for example \texttt{stressControlAdaptGain = 1.0e-2}.  For
zero-stress targets, such as uniaxial-stress tension, specify a nonzero
\texttt{stressControlStressFloor} based on the material stress scale or measured
stress noise.

\paragraph{Anisotropic plane-strain compression.}
For a plane-strain graphite block compressed in $Y$ with periodic relaxation in
$X$, use a hydrostatic hold, prescribe $F_{yy}$, and stress-control only the
free lateral direction:
\begin{lstlisting}[caption={Plane-strain anisotropic compression with X pressure control.}]
pfw["periodic"] = [True, False, False]
pfw["planeStrain"] = 1
pfw["prescribedBoundaryFTable"] = 1
pfw["prescribedFTable"] = 1
pfw["fTable"] = [
  [0.0,                 1.0, 1.0,         1.0],
  [hydrostaticHoldTime, 1.0, 1.0,         1.0],
  [stopTime,            1.0, finalFyy,    1.0],
]
pfw["stressControl"] = [1, 0, 0]
pfw["stressTable"] = [[0.0, -Pc, -Pc, -Pc], [stopTime, -Pc, -Pc, -Pc]]
pfw["stressControlAdaptiveSeek"] = 1
pfw["stressControlKp"] = pfw["stressControlKi"] = pfw["stressControlKd"] = 0.0
pfw["stressControlResponseStrain"] = 1.0e-2
pfw["stressControlFilterStrain"] = 1.0e-3
pfw["stressControlMaxSeekStrain"] = 5.0e-2
\end{lstlisting}

\paragraph{Initially sparse granular compaction.}
For an initially unjammed granular body, stress-control all periodic directions
and provide a physical reference strain rate because no prescribed F-table rate
exists.  Estimate the per-direction seek allowance from the initial solid volume
fraction $\phi_0$ and an expected jamming fraction $\phi_j$:
\begin{lstlisting}[caption={Sparse granular compaction to hydrostatic pressure.}]
phi0 = solidVolume / domainVolume
phiJam = 0.58
nSeek = 3
maxSeek = 1.2 * max(0.0, math.log(phiJam / phi0)) / nSeek + 0.02

pfw["periodic"] = [True, True, True]
pfw["prescribedFTable"] = 1
pfw["fTable"] = [[0.0, 1.0, 1.0, 1.0], [stopTime, 1.0, 1.0, 1.0]]
pfw["stressControl"] = [1, 1, 1]
pfw["stressTable"] = [[0.0, -Ptarget, -Ptarget, -Ptarget],
                      [stopTime, -Ptarget, -Ptarget, -Ptarget]]
pfw["stressControlAdaptiveSeek"] = 1
pfw["stressControlMinStrainRate"] = seekStrainRate
pfw["stressControlMaxSeekStrain"] = maxSeek
pfw["stressControlSeekRateRatio"] = 1.0
pfw["stressControlAdaptGain"] = 1.0e-2
\end{lstlisting}

\paragraph{Triaxial compression of a geomaterial.}
For a 3D RVE with axial $Y$ compression and constant lateral confinement, control
both lateral stresses.  Do not use this setup in plane strain unless $Z$ is free
to deform:
\begin{lstlisting}[caption={RVE triaxial compression with X/Z confining pressure.}]
pfw["planeStrain"] = 0
pfw["periodic"] = [True, True, True]
pfw["prescribedFTable"] = 1
pfw["fTable"] = [[0.0, 1.0, 1.0, 1.0],
                 [holdTime, 1.0, 1.0, 1.0],
                 [stopTime, 1.0, finalFyy, 1.0]]
pfw["stressControl"] = [1, 0, 1]
pfw["stressTable"] = [[0.0, -Pc, -Pc, -Pc], [stopTime, -Pc, -Pc, -Pc]]
pfw["stressControlAdaptiveSeek"] = 1
pfw["stressControlAdaptGain"] = 1.0e-2
pfw["stressControlMaxSeekStrain"] = 2.0e-2
\end{lstlisting}

\paragraph{Uniaxial-stress tension with hardening and damage.}
For axial tension with zero lateral stress, control the lateral stresses to zero
and set a finite stress floor.  Because the target is not compressive, seek mode
does not compact the specimen after failure; if lateral stress falls to zero,
the target is already satisfied:
\begin{lstlisting}[caption={Uniaxial-stress tension with zero lateral stress.}]
pfw["planeStrain"] = 0
pfw["periodic"] = [True, True, True]
pfw["prescribedFTable"] = 1
pfw["fTable"] = [[0.0, 1.0, 1.0, 1.0], [stopTime, 1.0, finalFyy, 1.0]]
pfw["stressControl"] = [1, 0, 1]
pfw["stressTable"] = [[0.0, 0.0, 0.0, 0.0], [stopTime, 0.0, 0.0, 0.0]]
pfw["stressControlAdaptiveSeek"] = 1
pfw["stressControlStressFloor"] = max(1.0e-4 * E, 1.0e-3 * sigmaY)
pfw["stressControlResponseStrain"] = 5.0e-3
pfw["stressControlFilterStrain"] = 5.0e-4
pfw["stressControlMaxSeekStrain"] = 0.0
\end{lstlisting}

Persistent wave warnings indicate that the controlled stress is dominated by
transient stress waves rather than a quasi-static mean.  Increase
\texttt{stressControlResponseStrain} and \texttt{stressControlFilterStrain},
reduce rate caps, or slow the F-table ramp.  If seek mode exceeds
\texttt{stressControlMaxSeekStrain}, the requested compressive target is not
reachable within the allowed closure; for granular compaction increase the
volume-fraction-based allowance only if it is physically justified.

\subsection{Reaction accumulation and reaction-history output}
\label{subsec:bc-reactions-internals}
\index{reactionHistory!internal algorithm}

When a moving or contact boundary changes a normal nodal velocity, the solver can accumulate a face reaction.  Only nodes owned by the local rank, identified by \texttt{gridGhostRank <= -1}, and with nodal mass above \texttt{smallMass} contribute.  The default reaction calculation is the mass times the velocity correction divided by the time step,
\begin{equation}
  R_f^{\mathrm{local}} \mathrel{+}= m_{I\alpha}\frac{v_{I\alpha,d}^{\mathrm{bc}}-v_{I\alpha,d}^{\mathrm{old}}}{\Delta t}.
  \label{eq:bc-reaction-momentum}
\end{equation}
If \texttt{useInternalForceAsFaceReaction} is enabled, the contribution is instead based on the internal force component,
\begin{equation}
  R_f^{\mathrm{local}} \mathrel{-}= f_{I\alpha,d}^{\mathrm{int}}.
  \label{eq:bc-reaction-internal-force}
\end{equation}
The six local face sums are reduced across MPI ranks with a sum reduction and stored in \texttt{m\_globalFaceReactions}.  If \texttt{reactionHistory} is enabled and the write schedule is satisfied, the solver appends a row to \path{reactionHistory.csv} containing time, domain deformation components, end-of-step box lengths, the six face reactions, and the current diagonal \texttt{domainL} values.  For \texttt{prescribedBoundaryFTable}, the reported lengths use the same exact finite ratios as the grid resize, so the reported $\bar F$ and domain dimensions remain mutually consistent.  This output is therefore a diagnostic of the constraints actually applied by \texttt{applyEssentialBCs}; it is not produced for pure outflow or inactive moving faces that do not change normal nodal velocities.

\subsection{FMPM-specific boundary consistency}
\label{subsec:bc-fmpm-internals}
\index{FMPM!boundary conditions}

The FMPM update stores an uncontacted grid velocity seed before material-contact corrections.  After \texttt{applyEssentialBCs} has imposed symmetry or moving constraints on \texttt{gridVelocity}, the solver copies only the constrained boundary components back into \texttt{gridUncontactedVelocity}.  Unconstrained tangential components are deliberately left material-contact-free so that the filtered contact update can use them as an uncorrected seed.

Higher-order FMPM velocity increments use a homogeneous version of the boundary constraints.  On boundary nodes, constrained components of the increment are set to zero.  On buffer nodes, the normal increment is reflected oddly,
\begin{equation}
  \Delta v_{G\alpha,d}=-\Delta v_{J(G)\alpha,d},
\end{equation}
and transverse increments are either reflected oddly when transverse velocities are prescribed or left unchanged when they are free.  This is the same incremental compatibility idea emphasized by Nairn \cite{nairn2026fmpm}: grid-based velocity constraints and lumped-mass contact operations are applied to each FMPM increment so that the final filtered velocity still satisfies the required constraints.  Therefore, the FMPM restriction should not be read as a general incompatibility with moving velocity boundaries or multi-field material contact.  Those paths use the incremental correction strategy.  The code-status caveat is specifically the \texttt{boundaryConditionTypes == Contact} wall/contact-boundary bctype, which is guarded in the reviewed branch and should be treated as untested until the corresponding verification case is added.

\subsection{Periodic boundaries as topology rather than face constraints}
\label{subsec:bc-periodic-internals}
\index{periodic boundaries}

Periodic boundaries are not represented by \texttt{boundaryConditionTypes}.  They are stored in the spatial partition and affect grid and particle topology.  During initialization, the solver shifts periodic ghost-node coordinates on end partitions by one global domain extent so that distance-based operations see the closest periodic image.  During the topology-preparation step, inactive ghost particle centers are wrapped across periodic boundaries for neighbor operations.  During repartitioning and cleanup, active particle centers are wrapped back into the global domain by a modulo operation,
\begin{equation}
  x_{p,i}\leftarrow \operatorname{mod}(x_{p,i}-X_i^{\min},L_i)+X_i^{\min}
  \qquad\text{for periodic direction } i.
  \label{eq:bc-periodic-wrap}
\end{equation}
Out-of-range deletion checks are skipped in periodic directions.  If \texttt{prescribedFTable} is active, the superposed velocity-gradient update is applied to particles only in periodic directions, which is the internal distinction between domain-periodic homogeneous deformation and boundary-driven deformation on moving faces.
